% Capability-Specific Degradation Patterns in Quantized Small Language Models
% Built for the ACL 2023 LaTeX template. See paper/README.md for how to obtain
% the style files (acl.sty, acl_natbib.bst) so this compiles.
%
% Workflow: results -> CSVs -> analysis -> tables/figures -> this file.
%   python scripts/aggregate_results.py ...
%   python scripts/analyze_results.py ...
%   python analysis/make_paper_tables.py ...
% then compile:  pdflatex main && bibtex main && pdflatex main && pdflatex main

\documentclass[11pt]{article}

\usepackage[]{acl}            % ACL 2023 style; provides \aclfinalcopy etc.
\usepackage{times}
\usepackage{latexsym}
\usepackage[T1]{fontenc}
\usepackage[utf8]{inputenc}
\usepackage{microtype}
\usepackage{booktabs}
\usepackage{graphicx}
\usepackage{amsmath}
\usepackage{xcolor}

\title{Capability-Specific Degradation Patterns in Quantized Small Language Models}

\author{Emil Rahimov \\ \texttt{emilrahimov25@gmail.com}}

\begin{document}
\maketitle

\begin{abstract}
Post-training quantization is the default way to deploy small language models
(1--4B parameters), yet its impact is usually summarized by a single aggregate
accuracy number that can hide severe failures in individual capabilities. We
evaluate seven open instruction-tuned small LLMs spanning five architecture
families (Qwen2.5, Llama-3.2, Gemma-2, Phi-3.5, and SmolLM2) at FP16 and 4-bit
(bitsandbytes NF4) across six capabilities: factual knowledge, commonsense,
mathematical reasoning, multilingual math, code generation, and instruction
following. Rather than reporting only mean accuracy, we construct a
\emph{capability-level degradation map} and ask whether capabilities degrade in a
consistent order across architectures. We find that degradation is strongly
capability-specific: multilingual mathematical reasoning and code generation are
the most fragile capabilities (relative drops up to $57\%$), while commonsense
reasoning is almost perfectly preserved ($\sim$2\%) everywhere. However, the
\emph{ordering} of degradation is only weakly consistent across architectures
(mean pairwise Kendall's $\tau = 0.29$): the most- and least-fragile capabilities
are shared, but the middle ordering is architecture-dependent. We further show
that aggregate scores mask these failures---several models lose more than a fifth
of a specific capability while their mean barely moves---and we document a
measurement pitfall in which a heavily chat-tuned model scores near zero on code
under raw-completion prompting yet degrades by $57\%$ once prompted correctly. Our
results argue that small-model quantization should be reported per capability,
not in aggregate.
\end{abstract}

\section{Introduction}

Small language models in the 1--4B parameter range are increasingly the
deployment-relevant artifact: they run on a single consumer GPU, on laptops, and
on edge devices. Post-training quantization---most commonly to 4 bits---is the
standard way to fit these models into tight memory budgets. The practical
question for a practitioner is simple: \emph{what do I lose when I quantize?}

The usual answer is an aggregate benchmark score, or an average over a benchmark
suite. We argue this is the wrong granularity. A model whose average accuracy
drops by three points may have lost almost nothing on commonsense reasoning while
losing a quarter of its code-generation or multilingual ability. For a deployed
system that distinction is the whole story: an aggregate number that looks
acceptable can conceal a capability that has effectively broken.

We therefore study quantization at the level of individual capabilities. We
evaluate seven instruction-tuned small models from five architecture families at
FP16 and 4-bit across six capabilities, and we organize the results as a
\emph{degradation map}: for every model and capability, the relative drop from
FP16 to 4-bit. We then ask whether different capabilities degrade in a
\emph{consistent order} across architectures---if they do, the pattern is
universal and predictable; if they do not, degradation is architecture-dependent
and must be measured per model.

Our contributions are:
\begin{enumerate}
\item A capability-level degradation map for seven open small LLMs (five families)
under 4-bit quantization across six capabilities, using a single controlled
evaluation pipeline (Section~\ref{sec:setup}).
\item Evidence that degradation is strongly \emph{capability-specific}:
multilingual math and code degrade most, commonsense least, with more than a
$25\times$ spread in relative degradation within a single model
(Section~\ref{sec:degradation}).
\item A cross-architecture consistency analysis (mean Kendall's $\tau = 0.29$)
showing that the degradation \emph{ordering} is only weakly universal: the
extremes are shared but the middle ranking is architecture-dependent
(Section~\ref{sec:consistency}).
\item A size-axis analysis showing that smaller models degrade more in magnitude,
and a demonstration that aggregate scores hide capability-specific collapse,
including a prompting-protocol pitfall on code evaluation
(Sections~\ref{sec:size}--\ref{sec:hidden}).
\end{enumerate}

\section{Related Work}

\paragraph{Unified compression benchmarks.} UniComp~\citep{unicomp2026} compares
pruning, quantization, and distillation across 40+ datasets and reports a
``knowledge bias'': knowledge-style tasks are preserved while reasoning,
multilingual, and instruction-following abilities degrade. Their focus is
\emph{comparing methods} on largely mid-to-large models; they do not isolate the
small-model regime or produce an explicit per-capability degradation ranking,
which is our focus.

\paragraph{Quantization and reasoning.} \citet{quantreasoning2025} provide a
math-specific analysis, reporting up to 32\% degradation on Llama-3 mathematical
reasoning under AWQ/GPTQ. We corroborate the fragility of mathematical reasoning
but situate it among five other capabilities and across architectures and model
sizes.

\paragraph{Multilingual quantization.} \citet{marchisio2024multilingual} show
that quantization disproportionately harms non-Latin-script languages and that
automatic metrics can underestimate the damage. Consistent with this, we find
multilingual math (MGSM) to be the single most fragile capability in our suite.

\paragraph{Coding and STEM.} A recent quantization benchmark~\citep{ijcai25quant}
finds coding and STEM tasks degrade most on MT-Bench, but is Llama-centric and
oriented toward comparing quantization methods rather than ranking capabilities
across architectures.

\paragraph{Mechanisms.} \citet{signaldegradation2026} give a mechanistic account
of why extreme (2-bit) quantization fails, identifying distinct failure modes.
Our study is complementary: we measure benchmark-level capability degradation at
the deployment-relevant 4-bit setting rather than analyzing internal mechanisms.

\paragraph{Positioning.} Relative to this work we (i) restrict to the understudied
1--4B small-model regime, (ii) compare five architecture families under one
identical pipeline, (iii) report an explicit capability degradation
\emph{ordering} and test its cross-architecture consistency, and (iv) add a size
axis within two families.

\section{Experimental Setup}
\label{sec:setup}

\paragraph{Models.} We evaluate seven instruction-tuned models. Five form the
cross-architecture set used for the consistency analysis: Qwen2.5-3B-Instruct,
Llama-3.2-3B-Instruct, Gemma-2-2B-it, Phi-3.5-mini-instruct (3.8B), and
SmolLM2-1.7B-Instruct. Two additional models form a size axis within existing
families: Qwen2.5-1.5B-Instruct and Llama-3.2-1B-Instruct. We use
instruction-tuned variants because instruction following (IFEval) requires them
and because they are the deployment-relevant artifact at this scale. All models
are $\leq 4$B parameters by design; Gemma-2 is loaded with eager attention so its
logit soft-capping is applied correctly.

\paragraph{Quantization.} We compare two precisions: FP16 (baseline) and 4-bit
weight quantization via bitsandbytes NF4~\citep{dettmers2022llmint8}. We focus on
the aggressive 4-bit regime because it is where capability-specific failures
emerge; 8-bit quantization was near-lossless in pilot runs and is omitted. We
note GPTQ~\citep{frantar2023gptq} and AWQ~\citep{lin2024awq} as alternative 4-bit
methods; we standardize on NF4 to hold the quantization method fixed across all
architectures, isolating architecture as the variable of interest.

\paragraph{Capabilities and benchmarks.} We map six capabilities to standard
benchmarks, evaluated with the LM Evaluation Harness~\citep{lmeval}: factual
knowledge (ARC-Challenge~\citep{clark2018arc}, 25-shot, acc\_norm); commonsense
(HellaSwag~\citep{zellers2019hellaswag}, 10-shot, acc\_norm); mathematical
reasoning (GSM8K~\citep{cobbe2021gsm8k}, 8-shot, flexible-extract exact match);
multilingual math (MGSM~\citep{shi2022mgsm}, 8-shot, exact match, macro-averaged
over the standard languages); code generation
(HumanEval~\citep{chen2021humaneval}, 0-shot, pass@1); and instruction following
(IFEval~\citep{zhou2023ifeval}, 0-shot, prompt-level strict accuracy). Generative
tasks (GSM8K, MGSM, IFEval) use the chat template; HumanEval uses raw completion
(the standard protocol); multiple-choice tasks use loglikelihood scoring. MGSM
uses its standard language set and does not include Azerbaijani or Turkish.

\paragraph{Evaluation pipeline.} All runs use a single controlled pipeline.
Generative tasks are executed with a high-throughput vLLM backend; multiple-choice
loglikelihood tasks use the Hugging Face Transformers backend. The same backend is
used for a given capability across all precisions and models, so any
FP16-vs-4-bit difference within a capability reflects quantization and not the
harness. We verified pipeline correctness by matching FP16 baselines against
published numbers (e.g., Qwen2.5-3B GSM8K and ARC). All $7\times2\times6=84$ runs
completed.

\paragraph{Analysis.} For each model and capability we compute the relative
degradation
\[
\Delta_{\text{rel}} = \frac{s_{\text{FP16}} - s_{\text{4bit}}}{s_{\text{FP16}}}\times 100\%.
\]
To avoid unstable ratios, we exclude any capability whose FP16 baseline is below
10 points from the relative-degradation analysis (these are reported as absolute
scores) and additionally report absolute drops as a sanity check. We rank
capabilities by $\Delta_{\text{rel}}$ within each model (rank 1 = most degraded)
and compare rankings across models with Kendall's $\tau$~\citep{kendall1938tau}.
A high mean $\tau$ indicates a universal degradation order; a low mean $\tau$
indicates architecture-dependence. We separately compute $\tau$ between the small
and large member of each family (the size axis).

\section{Results}

\subsection{Capability-Specific Degradation}
\label{sec:degradation}

Table~\ref{tab:degradation} reports the relative degradation map and
Table~\ref{tab:absolute} the underlying absolute scores. The headline observation
is the very large spread \emph{within} a single model. For Qwen2.5-1.5B, relative
degradation ranges from $2.6\%$ (HellaSwag) to $22.9\%$ (MGSM) and $22.6\%$
(HumanEval)---an order-of-magnitude difference between the most- and
least-affected capability. Aggregating these into one number would report a
moderate drop and erase the fact that two capabilities lost roughly a quarter of
their performance.

Two capabilities stand out as consistently fragile. \textbf{Multilingual math
(MGSM)} is the most degraded capability for five of the six models for which it
is measurable, with relative drops of $12$--$23\%$. \textbf{Code generation
(HumanEval)} is the most volatile: catastrophic for some models (Qwen2.5-1.5B
$22.6\%$, Phi-3.5 $21.0\%$, and SmolLM2 $57\%$; see Section~\ref{sec:hidden}) yet
statistically unchanged for others (Llama-3.2-3B, Gemma-2). At the other extreme,
\textbf{commonsense (HellaSwag)} is essentially preserved everywhere
($1.7$--$5.1\%$), and factual knowledge (ARC) is largely robust. Small negative
values (4-bit marginally above FP16) occur for a few code/math cells and are
within evaluation noise. Figure~\ref{fig:heatmap} visualizes the map.

\begin{table*}[t]
\centering
\small
\begin{tabular}{lrrrrrr}
\toprule
Model (4-bit vs FP16) & GSM8K & ARC & HellaSwag & HumanEval & IFEval & MGSM \\
\midrule
Qwen2.5-3B   & $-2.6$ & $3.4$  & $2.6$ & $9.2$   & $1.9$  & $19.0$ \\
Qwen2.5-1.5B & $10.6$ & $9.3$  & $2.6$ & $22.6$  & $11.8$ & $22.9$ \\
Llama-3.2-3B & $4.0$  & $3.1$  & $1.8$ & $-6.1$  & $-0.5$ & $19.1$ \\
Llama-3.2-1B & $16.7$ & $2.7$  & $5.1$ & $12.3$  & $4.5$  & --     \\
Gemma-2-2B   & $4.1$  & $2.4$  & $1.8$ & $-2.6$  & $1.0$  & $12.2$ \\
Phi-3.5-mini & $4.3$  & $-1.7$ & $1.7$ & $21.0$  & $1.1$  & $19.0$ \\
SmolLM2-1.7B & $18.1$ & $2.7$  & $2.4$ & $57.2^{\dagger}$ & $-1.2$ & --   \\
\bottomrule
\end{tabular}
\caption{Relative degradation $\Delta_{\text{rel}}$ (\%) from FP16 to 4-bit
(higher = worse). ``--'' marks capabilities excluded for an FP16 baseline below
10 points (MGSM for the two smallest models). $\dagger$: SmolLM2 HumanEval under
the chat-template protocol (Section~\ref{sec:hidden}); under the common
raw-completion protocol its baseline is near zero and the cell is excluded.}
\label{tab:degradation}
\end{table*}

\begin{table*}[t]
\centering
\small
\begin{tabular}{llrrrrrr}
\toprule
Model & Prec. & GSM8K & ARC & HellaSwag & HumanEval & IFEval & MGSM \\
\midrule
Qwen2.5-3B    & FP16 & 63.4 & 60.7 & 75.3 & 53.1 & 59.3 & 24.9 \\
              & 4-bit & 65.1 & 58.6 & 73.3 & 48.2 & 58.2 & 20.2 \\
Qwen2.5-1.5B  & FP16 & 61.7 & 54.4 & 67.9 & 37.8 & 40.9 & 10.3 \\
              & 4-bit & 55.2 & 49.3 & 66.1 & 29.3 & 36.0 & 8.0 \\
Llama-3.2-3B  & FP16 & 77.3 & 52.8 & 73.6 & 50.0 & 69.3 & 31.0 \\
              & 4-bit & 74.2 & 51.2 & 72.3 & 53.1 & 69.7 & 25.1 \\
Llama-3.2-1B  & FP16 & 44.1 & 41.8 & 61.1 & 34.8 & 49.9 & 6.4 \\
              & 4-bit & 36.8 & 40.7 & 57.9 & 30.5 & 47.7 & 6.0 \\
Gemma-2-2B    & FP16 & 54.1 & 57.9 & 71.5 & 23.8 & 55.6 & 37.6 \\
              & 4-bit & 51.9 & 56.6 & 70.2 & 24.4 & 55.1 & 33.1 \\
Phi-3.5-mini  & FP16 & 78.6 & 64.8 & 79.0 & 61.0 & 49.9 & 33.0 \\
              & 4-bit & 75.3 & 65.9 & 77.7 & 48.2 & 49.4 & 26.8 \\
SmolLM2-1.7B  & FP16 & 49.1 & 53.2 & 72.6 & $34.2^{\dagger}$ & 48.2 & 3.9 \\
              & 4-bit & 40.3 & 51.8 & 70.8 & $14.6^{\dagger}$ & 48.8 & 3.2 \\
\bottomrule
\end{tabular}
\caption{Absolute scores (\%) per model and precision. $\dagger$: SmolLM2
HumanEval under the chat-template protocol; under raw completion both values are
near zero (see Section~\ref{sec:hidden}).}
\label{tab:absolute}
\end{table*}

\begin{figure}[t]
\centering
\includegraphics[width=\columnwidth]{figures/degradation_heatmap.png}
\caption{Relative degradation (\%) per model (rows) and capability (columns).
Multilingual math and code are the hot columns; commonsense is uniformly cool.}
\label{fig:heatmap}
\end{figure}

\subsection{Cross-Architecture Consistency}
\label{sec:consistency}

We rank the six capabilities by degradation within each model and compare
rankings across all $\binom{7}{2}=21$ model pairs with Kendall's $\tau$. The mean
pairwise $\tau$ is $\mathbf{0.29}$, which we interpret as \emph{weak} consistency:
the degradation order is \emph{not} universal across architectures. The pattern is
better described as partially consistent. The extremes are shared---multilingual
math is reliably near the top of the ranking and commonsense reliably near the
bottom---but the middle of the ranking (knowledge, math, code, instruction
following) reorders from architecture to architecture, which depresses the mean
$\tau$. Pairwise agreement is correspondingly heterogeneous: some pairs rank
capabilities identically (e.g., Gemma-2 vs.\ Llama-3.2-3B, $\tau=1.0$) while
others are uncorrelated (near $\tau=0$). Figure~\ref{fig:ranking} shows the
rankings side by side.

The practical implication is that one cannot reliably predict a new
architecture's full degradation profile from another's. A deployment decision
should measure the specific model rather than assume a universal ordering. This is
itself actionable: it argues against ``rules of thumb'' for which capabilities
survive quantization.

\begin{figure}[t]
\centering
\includegraphics[width=\columnwidth]{figures/ranking_comparison.png}
\caption{Degradation rankings per model (1 = most degraded). Agreement across
columns would indicate a shared ordering; the visible reordering in the middle
ranks reflects the low mean Kendall's $\tau$.}
\label{fig:ranking}
\end{figure}

\subsection{Size Axis}
\label{sec:size}

Within the two families evaluated at two sizes, the smaller model degrades
\emph{more in magnitude}. Qwen2.5-1.5B degrades more than Qwen2.5-3B on five of
six capabilities (e.g., GSM8K $10.6\%$ vs.\ $-2.6\%$, HumanEval $22.6\%$ vs.\
$9.2\%$); Llama-3.2-1B degrades more than Llama-3.2-3B on math and code (GSM8K
$16.7\%$ vs.\ $4.0\%$). The degradation \emph{ordering}, however, does not
strongly persist across sizes: the small-vs-large Kendall's $\tau$ is $0.47$ for
Qwen2.5 and $0.00$ for Llama-3.2. Quantization thus hurts smaller models more,
but not in a fixed capability order even within a family.

\subsection{Aggregate Scores Hide Capability Failures}
\label{sec:hidden}

The degradation map exposes failures that an aggregate would conceal. The clearest
case is code generation under quantization for small chat-tuned models.

Under the standard raw-completion HumanEval protocol used for all models,
SmolLM2-1.7B scores $0.6\%$ at FP16 and $0.0\%$ at 4-bit---numbers that, taken at
face value, suggest the model simply cannot code. Inspection of its generations
shows the cause is a \emph{prompting-protocol artifact}: SmolLM2 is heavily
chat-tuned and, given a bare function signature to continue, emits prose and
markdown rather than code. The six other models continue raw code normally. When
SmolLM2 is instead prompted with its chat template and the emitted function is
extracted and executed, it scores $34.2\%$ at FP16 and $14.6\%$ at 4-bit---a
\textbf{relative degradation of $57\%$, the largest code degradation in our
study}. The raw-completion protocol had hidden a near-total collapse of a real
capability.

We draw two lessons. First, capability-specific evaluation can be sensitive to
prompting protocol, and a single near-zero score should be audited before being
read as inability. Second, and more importantly for our thesis, the underlying
phenomenon is severe: a 1.7B model loses more than half of a genuine
code-generation ability under 4-bit quantization, which no aggregate score in our
suite surfaces. We report SmolLM2's chat-protocol code result separately (the
$\dagger$ entries in Tables~\ref{tab:degradation} and~\ref{tab:absolute}) because
it uses a different prompt format from the other models; the cross-model
consistency analysis (Section~\ref{sec:consistency}) uses only the common
protocol.

\section{Discussion}

Three themes emerge. \textbf{(1) Capability, not aggregate.} Across every model,
quantization damage is concentrated in a minority of capabilities. Reporting a
single accuracy number is not just lossy but actively misleading when the lost
capability is the one a deployment depends on (e.g., multilingual or code).
\textbf{(2) Fragility is partly intrinsic to the capability.} Multilingual math
and code are fragile across architectures while commonsense is robust, suggesting
the brittleness is tied to what the capability demands---long multi-step
generation and exact-token correctness---more than to the architecture.
\textbf{(3) But ordering is architecture-dependent.} The weak mean $\tau$ shows
that beyond the shared extremes, which capability breaks next depends on the
model, so practitioners cannot port a degradation profile across families. The
deployment recommendation that follows is concrete: before shipping a quantized
small model, re-test its multilingual and code capabilities specifically, and do
not rely on an aggregate score or on another model's profile.

\section*{Limitations}

Our study covers seven models, one quantization method (NF4) at one bit-width
(4-bit), and one benchmark per capability; broader coverage (AWQ/GPTQ, additional
benchmarks per capability, more families) would strengthen the consistency
estimates, whose per-pair $p$-values are limited by having only six capabilities.
We report single evaluation runs; generative benchmarks (GSM8K, HumanEval, MGSM)
carry run-to-run variance, so small ranking differences should be read with the
reported standard errors in mind. MGSM uses the standard language set and is
excluded for the two smallest models, where the FP16 baseline is too low to form
a stable ratio; degradation there is reported in absolute terms only. The
Kendall's-$\tau$ analysis ranks six capabilities, so individual pairwise
correlations are noisy and we emphasize the mean. Finally, as
Section~\ref{sec:hidden} shows, code evaluation is protocol-sensitive for
chat-tuned models; we mitigated this by auditing and re-prompting, but a fully
uniform instruction-style protocol across all models is left to future work.

\section{Conclusion}

We presented a capability-level degradation map for seven small instruction-tuned
LLMs under 4-bit quantization. Degradation is strongly capability-specific:
multilingual math and code are most fragile (up to $57\%$ relative loss) while
commonsense is nearly untouched. The degradation \emph{ordering} is only weakly
consistent across architectures (mean Kendall's $\tau=0.29$), so the extremes are
predictable but the full profile is architecture-dependent. Smaller models
degrade more in magnitude. Because aggregate accuracy hides these
capability-specific failures, we recommend that small-model quantization be
reported and audited per capability.

\bibliography{references}
\bibliographystyle{acl_natbib}

\end{document}
